7. (1.5 pts) Considere \( f \in \operatorname{End}\left(M_{2 \times 2}(\mathbb{R})\right) \) dado por \( f(A)=A+A^{t} \). Encuentre:\\
- Polinomio característico.\\
- Valores y vectores propios.\\
- Espacios propios asociados a cada valor propio.\\

\begin{solution}
Consideremos la base canónica de \( M_{2 \times 2}(\mathbb{R}) \):
\[
   E_{11} = \begin{pmatrix} 1 & 0 \\ 0 & 0 \end{pmatrix}, \quad E_{12} = \begin{pmatrix} 0 & 1 \\ 0 & 0 \end{pmatrix}, \quad E_{21} = \begin{pmatrix} 0 & 0 \\ 1 & 0 \end{pmatrix}, \quad E_{22} = \begin{pmatrix} 0 & 0 \\ 0 & 1 \end{pmatrix}
   \]
   Calculemos \( f \) en cada uno:
   \[
   f(E_{11}) = \begin{pmatrix} 1 & 0 \\ 0 & 0 \end{pmatrix} + \begin{pmatrix} 1 & 0 \\ 0 & 0 \end{pmatrix} = \begin{pmatrix} 2 & 0 \\ 0 & 0 \end{pmatrix}
   \]
\[
   f(E_{12}) = \begin{pmatrix} 0 & 1 \\ 0 & 0 \end{pmatrix} + \begin{pmatrix} 0 & 0 \\ 1 & 0 \end{pmatrix} = \begin{pmatrix} 0 & 1 \\ 1 & 0 \end{pmatrix}
   \]
\[
   f(E_{21}) = \begin{pmatrix} 0 & 0 \\ 1 & 0 \end{pmatrix} + \begin{pmatrix} 0 & 1 \\ 0 & 0 \end{pmatrix} = \begin{pmatrix} 0 & 1 \\ 1 & 0 \end{pmatrix}
   \]
\[
   f(E_{22}) = \begin{pmatrix} 0 & 0 \\ 0 & 1 \end{pmatrix} + \begin{pmatrix} 0 & 0 \\ 0 & 1 \end{pmatrix} = \begin{pmatrix} 0 & 0 \\ 0 & 2 \end{pmatrix}
   \]
   La matriz de f en la base canónica es:
\[
   \begin{pmatrix}
   2 & 0 & 0 & 0 \\
   0 & 0 & 1 & 0 \\
   0 & 1 & 0 & 0 \\
   0 & 0 & 0 & 2
   \end{pmatrix}
   \]
   El polinomio característico es \(\det(f - \lambda I)\):
\[
   \begin{vmatrix}
   2 - \lambda & 0 & 0 & 0 \\
   0 & -\lambda & 1 & 0 \\
   0 & 1 & -\lambda & 0 \\
   0 & 0 & 0 & 2 - \lambda
   \end{vmatrix}
   = (2 - \lambda)^2 (\lambda^2 - 1)
   \]
Por lo tanto, el polinomio característico es:
\[
   (\lambda - 2)^2 (\lambda^2 - 1)
   \]
Lo dejo en esta forma pues es directo calcular las raíces y su multiplicidad, algo que nos servirá para el cálculo de valores y vectores propios. \\
Para los valores propios igualamos a cero nuestro polinomio característico.
\[
   (\lambda - 2)^2 (\lambda^2 - 1) = 0
   \]
Entonces algún miembro debe ser cero, si $(\lambda - 2)^2=0$ entonces $\lambda_{1,2}=2$. Por otro lado si $(\lambda^2 - 1)=(\lambda-1)(\lambda+1)=0$, usando el mismo principio indica que ya sea el primer o el segundo término, debe de ser cero. Se sigue que $\lambda_{1,2}=0, \lambda_3=1, \lambda_4=-1$. \\

Ahora obténganse los Vectores Propios y Espacios Propios Asociados. \\
1,2)   Para $\lambda=\lambda_{1,2}$: \\
Resolvemos \((f - 2I)X = 0\).
\[
     \begin{pmatrix}
     0 & 0 & 0 & 0 \\
     0 & -2 & 1 & 0 \\
     0 & 1 & -2 & 0 \\
     0 & 0 & 0 & 0
     \end{pmatrix}
     \begin{pmatrix}
     x_1 \\
     x_2 \\
     x_3 \\
     x_4
     \end{pmatrix}
     = 0
     \]
La solución general es \( x_1, x_4 \) libres, \( x_2 = \frac{1}{2}x_3 \), y \( x_3 = 0 \). \\
Espacio propio: \(\operatorname{span}\left\{\begin{pmatrix} 1 \\ 0 \\ 0 \\ 0 \end{pmatrix}, \begin{pmatrix} 0 \\ 0 \\ 0 \\ 1 \end{pmatrix}\right\}\). \\

3)   Para $\lambda=\lambda_{3}$: \\
Resolvemos \((f - I)X = 0\).
\[
     \begin{pmatrix}
     1 & 0 & 0 & 0 \\
     0 & -1 & 1 & 0 \\
     0 & 1 & -1 & 0 \\
     0 & 0 & 0 & 1
     \end{pmatrix}
     \begin{pmatrix}
     x_1 \\
     x_2 \\
     x_3 \\
     x_4
     \end{pmatrix}
     = 0
     \]
La solución es \( x_1 = 0, x_4 = 0, x_2 = x_3 \). \\
Espacio propio: \(\operatorname{span}\left\{\begin{pmatrix} 0 \\ 1 \\ 1 \\ 0 \end{pmatrix}\right\}\). \\

3)   Para $\lambda=\lambda_{4}$: \\
Resolvemos \((f + I)X = 0\).
\[
     \begin{pmatrix}
     3 & 0 & 0 & 0 \\
     0 & 1 & 1 & 0 \\
     0 & 1 & 1 & 0 \\
     0 & 0 & 0 & 3
     \end{pmatrix}
     \begin{pmatrix}
     x_1 \\
     x_2 \\
     x_3 \\
     x_4
     \end{pmatrix}
     = 0
     \]
La solución es \( x_1 = 0, x_4 = 0, x_2 = -x_3 \). \\
Espacio propio: \(\operatorname{span}\left\{\begin{pmatrix} 0 \\ 1 \\ -1 \\ 0 \end{pmatrix}\right\}\).
\end{solution}